% --- Archon physics formalization source begin ---
% archon:physics
% archon:covers PhyXMiniProblems/problem_phyx_mini_0979.lean
% archon:source-report reports/phyx_mini/problem_phyx_mini_0979.source.json

\chapter{Physics problem phyx\_mini\_0979}
\label{ch:PhyXMiniProblems_problem_phyx_mini_0979}

\paragraph{Problem source.}
\#\# Physical scenario

A bar with mass \textbackslash{}( M = 1.20\textbackslash{} \textbackslash{}text\{kg\} \textbackslash{}) and resistance \textbackslash{}( R = 0.500\textbackslash{} \textbackslash{}Omega \textbackslash{}) slides without friction on a horizontal U-shaped rail with width \textbackslash{}( W = 40.0\textbackslash{} \textbackslash{}text\{cm\} \textbackslash{}) and negligible resistance. The bar is attached to a spring with spring constant \textbackslash{}( k = 90.0\textbackslash{} \textbackslash{}text\{N/m\} \textbackslash{}). A constant magnetic field with magnitude \textbackslash{}( 1.00\textbackslash{} \textbackslash{}text\{T\} \textbackslash{}) points into the plane everywhere in the vicinity. At time \textbackslash{}( t = 0 \textbackslash{}) the bar is stretched beyond its equilibrium position by an amount \textbackslash{}( x = 10.0\textbackslash{} \textbackslash{}text\{cm\} \textbackslash{}) and released from rest.

\#\# Question

What is the amplitude of the motion at time \textbackslash{}( t = 5.00\textbackslash{} \textbackslash{}text\{s\} \textbackslash{})?

\#\# Answer choices

- A: A: 10.0cm

- B: B: 5.13cm

- C: C: 0cm

- D: D: 1.2cm

\#\# Figure caption (auxiliary; use the image as primary evidence)

The image depicts a U-shaped orange wire loop placed vertically in a magnetic field represented by blue Xs, indicating that the field (\textbackslash{}( \textbackslash{}vec\{B\} \textbackslash{})) is directed into the page. 

- The loop has a movable vertical side marked with \textbackslash{}( M, R \textbackslash{}).
- This movable side is attached to a spring on the left, which is connected to a fixed wall.
- The spring is labeled with \textbackslash{}( k \textbackslash{}).
- The horizontal sides of the loop are positioned such that the spring compresses or stretches along the horizontal axis \textbackslash{}( x \textbackslash{}).
- The loop's height or the distance between its horizontal sides is labeled \textbackslash{}( W \textbackslash{}).

This setup likely illustrates electromagnetic induction or force on a conductor in a magnetic field.

\#\# Dataset metadata

Electromagnetic Induction; Physical Model Grounding Reasoning, Multi-Formula Reasoning

\paragraph{Recorded answer/context.}
B: B: 5.13cm

\paragraph{Figure/image path.}
/root/proposal\_for\_physic/science-mango/phyx\_mini\_run/phyx\_data/test\_image/979.png

\paragraph{Formalization target.}
create a compiling Lean file with sorry bodies at `PhyXMiniProblems/problem\_phyx\_mini\_0979.lean`. The Lean declarations must preserve the physical quantities, dimensions or dimensional roles, figure labels, governing-law hypotheses, and final relation expressed by this problem.
Use Mathlib/Physlib names found through LeanExplore where available. If a physics API is missing, introduce faithful local abstractions rather than scalar placeholder aliases.

\begin{theorem}[Physics formalization target]
\label{thm:physics:phyx_mini_0979:target}
\lean{PhyXMiniProblems.ProblemPhyXMini0979.amplitude_at_five_seconds}
\uses{def:physics:phyx-mini-0979:phyxminiproblems-problemphyxmini0979-electromagneticspringsetup, def:physics:phyx-mini-0979:phyxminiproblems-problemphyxmini0979-matchesprimaryelectromagneticspringfigure, def:physics:phyx-mini-0979:phyxminiproblems-problemphyxmini0979-matcheselectromagneticspringscenario, def:physics:phyx-mini-0979:phyxminiproblems-problemphyxmini0979-electromagneticspringproblemdata, def:physics:phyx-mini-0979:phyxminiproblems-problemphyxmini0979-electromagneticspringlaws, def:physics:phyx-mini-0979:phyxminiproblems-problemphyxmini0979-roundstohundredthcentimeter}
The assigned autoformalize agent should translate this physics problem into Lean declarations in the covered file, with theorem and lemma proof bodies written as `by sorry`.
\end{theorem}
\begin{proof}
This is an autoformalization task, not a proof task. Produce statements that can later be proved by the physics prover without weakening the physical model.
\end{proof}

% --- Archon declaration topology begin ---
\section{Declaration topology}
\begin{definition}[magnetic Flux Density Dimension]
  \label{def:physics:phyx-mini-0979:phyxminiproblems-problemphyxmini0979-magneticfluxdensitydimension}
  \lean{PhyXMiniProblems.ProblemPhyXMini0979.magneticFluxDensityDimension}
  Magnetic flux density has dimension M T⁻¹ C⁻¹ (tesla)
\end{definition}
\begin{proof}
  This is the declared model definition of magnetic Flux Density Dimension.
\end{proof}
\begin{definition}[electrical Resistance Dimension]
  \label{def:physics:phyx-mini-0979:phyxminiproblems-problemphyxmini0979-electricalresistancedimension}
  \lean{PhyXMiniProblems.ProblemPhyXMini0979.electricalResistanceDimension}
  Electrical resistance has dimension M L² T⁻¹ C⁻² (ohm)
\end{definition}
\begin{proof}
  This is the declared model definition of electrical Resistance Dimension.
\end{proof}
\begin{definition}[spring Constant Dimension]
  \label{def:physics:phyx-mini-0979:phyxminiproblems-problemphyxmini0979-springconstantdimension}
  \lean{PhyXMiniProblems.ProblemPhyXMini0979.springConstantDimension}
  A spring constant has dimension M T⁻² (newton per metre)
\end{definition}
\begin{proof}
  This is the declared model definition of spring Constant Dimension.
\end{proof}
\begin{definition}[electromotive Force Dimension]
  \label{def:physics:phyx-mini-0979:phyxminiproblems-problemphyxmini0979-electromotiveforcedimension}
  \lean{PhyXMiniProblems.ProblemPhyXMini0979.electromotiveForceDimension}
  Electromotive force has dimension M L² T⁻² C⁻¹ (volt)
\end{definition}
\begin{proof}
  This is the declared model definition of electromotive Force Dimension.
\end{proof}
\begin{definition}[electric Current Dimension]
  \label{def:physics:phyx-mini-0979:phyxminiproblems-problemphyxmini0979-electriccurrentdimension}
  \lean{PhyXMiniProblems.ProblemPhyXMini0979.electricCurrentDimension}
  Electric current has dimension C T⁻¹ (ampere)
\end{definition}
\begin{proof}
  This is the declared model definition of electric Current Dimension.
\end{proof}
\begin{definition}[force Dimension]
  \label{def:physics:phyx-mini-0979:phyxminiproblems-problemphyxmini0979-forcedimension}
  \lean{PhyXMiniProblems.ProblemPhyXMini0979.forceDimension}
  Force has dimension M L T⁻² (newton)
\end{definition}
\begin{proof}
  This is the declared model definition of force Dimension.
\end{proof}
\begin{definition}[damping Coefficient Dimension]
  \label{def:physics:phyx-mini-0979:phyxminiproblems-problemphyxmini0979-dampingcoefficientdimension}
  \lean{PhyXMiniProblems.ProblemPhyXMini0979.dampingCoefficientDimension}
  A viscous damping coefficient has dimension M T⁻¹ (kg/s)
\end{definition}
\begin{proof}
  This is the declared model definition of damping Coefficient Dimension.
\end{proof}
\begin{definition}[Mass Magnitude]
  \label{def:physics:phyx-mini-0979:phyxminiproblems-problemphyxmini0979-massmagnitude}
  \lean{PhyXMiniProblems.ProblemPhyXMini0979.MassMagnitude}
  A nonnegative, unit-independent physical mass
\end{definition}
\begin{proof}
  This is the declared model definition of Mass Magnitude.
\end{proof}
\begin{definition}[Length Magnitude]
  \label{def:physics:phyx-mini-0979:phyxminiproblems-problemphyxmini0979-lengthmagnitude}
  \lean{PhyXMiniProblems.ProblemPhyXMini0979.LengthMagnitude}
  A nonnegative, unit-independent physical length
\end{definition}
\begin{proof}
  This is the declared model definition of Length Magnitude.
\end{proof}
\begin{definition}[Time Magnitude]
  \label{def:physics:phyx-mini-0979:phyxminiproblems-problemphyxmini0979-timemagnitude}
  \lean{PhyXMiniProblems.ProblemPhyXMini0979.TimeMagnitude}
  A nonnegative, unit-independent elapsed or clock time
\end{definition}
\begin{proof}
  This is the declared model definition of Time Magnitude.
\end{proof}
\begin{definition}[Speed Magnitude]
  \label{def:physics:phyx-mini-0979:phyxminiproblems-problemphyxmini0979-speedmagnitude}
  \lean{PhyXMiniProblems.ProblemPhyXMini0979.SpeedMagnitude}
  A nonnegative, unit-independent speed magnitude
\end{definition}
\begin{proof}
  This is the declared model definition of Speed Magnitude.
\end{proof}
\begin{definition}[Magnetic Flux Density Magnitude]
  \label{def:physics:phyx-mini-0979:phyxminiproblems-problemphyxmini0979-magneticfluxdensitymagnitude}
  \lean{PhyXMiniProblems.ProblemPhyXMini0979.MagneticFluxDensityMagnitude}
  \uses{def:physics:phyx-mini-0979:phyxminiproblems-problemphyxmini0979-magneticfluxdensitydimension}
  A nonnegative, unit-independent magnetic-flux-density magnitude
\end{definition}
\begin{proof}
  This is the declared model definition of Magnetic Flux Density Magnitude.
\end{proof}
\begin{definition}[Resistance Magnitude]
  \label{def:physics:phyx-mini-0979:phyxminiproblems-problemphyxmini0979-resistancemagnitude}
  \lean{PhyXMiniProblems.ProblemPhyXMini0979.ResistanceMagnitude}
  \uses{def:physics:phyx-mini-0979:phyxminiproblems-problemphyxmini0979-electricalresistancedimension}
  A nonnegative, unit-independent electrical resistance
\end{definition}
\begin{proof}
  This is the declared model definition of Resistance Magnitude.
\end{proof}
\begin{definition}[Spring Constant Magnitude]
  \label{def:physics:phyx-mini-0979:phyxminiproblems-problemphyxmini0979-springconstantmagnitude}
  \lean{PhyXMiniProblems.ProblemPhyXMini0979.SpringConstantMagnitude}
  \uses{def:physics:phyx-mini-0979:phyxminiproblems-problemphyxmini0979-springconstantdimension}
  A nonnegative, unit-independent spring constant
\end{definition}
\begin{proof}
  This is the declared model definition of Spring Constant Magnitude.
\end{proof}
\begin{definition}[Electromotive Force Magnitude]
  \label{def:physics:phyx-mini-0979:phyxminiproblems-problemphyxmini0979-electromotiveforcemagnitude}
  \lean{PhyXMiniProblems.ProblemPhyXMini0979.ElectromotiveForceMagnitude}
  \uses{def:physics:phyx-mini-0979:phyxminiproblems-problemphyxmini0979-electromotiveforcedimension}
  A nonnegative, unit-independent electromotive-force magnitude
\end{definition}
\begin{proof}
  This is the declared model definition of Electromotive Force Magnitude.
\end{proof}
\begin{definition}[Electric Current Magnitude]
  \label{def:physics:phyx-mini-0979:phyxminiproblems-problemphyxmini0979-electriccurrentmagnitude}
  \lean{PhyXMiniProblems.ProblemPhyXMini0979.ElectricCurrentMagnitude}
  \uses{def:physics:phyx-mini-0979:phyxminiproblems-problemphyxmini0979-electriccurrentdimension}
  A nonnegative, unit-independent electric-current magnitude
\end{definition}
\begin{proof}
  This is the declared model definition of Electric Current Magnitude.
\end{proof}
\begin{definition}[Force Magnitude]
  \label{def:physics:phyx-mini-0979:phyxminiproblems-problemphyxmini0979-forcemagnitude}
  \lean{PhyXMiniProblems.ProblemPhyXMini0979.ForceMagnitude}
  \uses{def:physics:phyx-mini-0979:phyxminiproblems-problemphyxmini0979-forcedimension}
  A nonnegative, unit-independent magnetic-force magnitude
\end{definition}
\begin{proof}
  This is the declared model definition of Force Magnitude.
\end{proof}
\begin{definition}[Damping Coefficient Magnitude]
  \label{def:physics:phyx-mini-0979:phyxminiproblems-problemphyxmini0979-dampingcoefficientmagnitude}
  \lean{PhyXMiniProblems.ProblemPhyXMini0979.DampingCoefficientMagnitude}
  \uses{def:physics:phyx-mini-0979:phyxminiproblems-problemphyxmini0979-dampingcoefficientdimension}
  A nonnegative, unit-independent mechanical damping coefficient
\end{definition}
\begin{proof}
  This is the declared model definition of Damping Coefficient Magnitude.
\end{proof}
\begin{definition}[mass In Kilograms]
  \label{def:physics:phyx-mini-0979:phyxminiproblems-problemphyxmini0979-massinkilograms}
  \lean{PhyXMiniProblems.ProblemPhyXMini0979.massInKilograms}
  \uses{def:physics:phyx-mini-0979:phyxminiproblems-problemphyxmini0979-massmagnitude}
  Read a mass in kilograms
\end{definition}
\begin{proof}
  This is the declared model definition of mass In Kilograms.
\end{proof}
\begin{definition}[length In Meters]
  \label{def:physics:phyx-mini-0979:phyxminiproblems-problemphyxmini0979-lengthinmeters}
  \lean{PhyXMiniProblems.ProblemPhyXMini0979.lengthInMeters}
  \uses{def:physics:phyx-mini-0979:phyxminiproblems-problemphyxmini0979-lengthmagnitude}
  Read a length in metres
\end{definition}
\begin{proof}
  This is the declared model definition of length In Meters.
\end{proof}
\begin{definition}[time In Seconds]
  \label{def:physics:phyx-mini-0979:phyxminiproblems-problemphyxmini0979-timeinseconds}
  \lean{PhyXMiniProblems.ProblemPhyXMini0979.timeInSeconds}
  \uses{def:physics:phyx-mini-0979:phyxminiproblems-problemphyxmini0979-timemagnitude}
  Read a time in seconds
\end{definition}
\begin{proof}
  This is the declared model definition of time In Seconds.
\end{proof}
\begin{definition}[speed In Meters Per Second]
  \label{def:physics:phyx-mini-0979:phyxminiproblems-problemphyxmini0979-speedinmeterspersecond}
  \lean{PhyXMiniProblems.ProblemPhyXMini0979.speedInMetersPerSecond}
  \uses{def:physics:phyx-mini-0979:phyxminiproblems-problemphyxmini0979-speedmagnitude}
  Read a speed in metres per second
\end{definition}
\begin{proof}
  This is the declared model definition of speed In Meters Per Second.
\end{proof}
\begin{definition}[magnetic Flux Density In Teslas]
  \label{def:physics:phyx-mini-0979:phyxminiproblems-problemphyxmini0979-magneticfluxdensityinteslas}
  \lean{PhyXMiniProblems.ProblemPhyXMini0979.magneticFluxDensityInTeslas}
  \uses{def:physics:phyx-mini-0979:phyxminiproblems-problemphyxmini0979-magneticfluxdensitymagnitude}
  Read magnetic flux density in teslas
\end{definition}
\begin{proof}
  This is the declared model definition of magnetic Flux Density In Teslas.
\end{proof}
\begin{definition}[resistance In Ohms]
  \label{def:physics:phyx-mini-0979:phyxminiproblems-problemphyxmini0979-resistanceinohms}
  \lean{PhyXMiniProblems.ProblemPhyXMini0979.resistanceInOhms}
  \uses{def:physics:phyx-mini-0979:phyxminiproblems-problemphyxmini0979-resistancemagnitude}
  Read an electrical resistance in ohms
\end{definition}
\begin{proof}
  This is the declared model definition of resistance In Ohms.
\end{proof}
\begin{definition}[spring Constant In Newtons Per Meter]
  \label{def:physics:phyx-mini-0979:phyxminiproblems-problemphyxmini0979-springconstantinnewtonspermeter}
  \lean{PhyXMiniProblems.ProblemPhyXMini0979.springConstantInNewtonsPerMeter}
  \uses{def:physics:phyx-mini-0979:phyxminiproblems-problemphyxmini0979-springconstantmagnitude}
  Read a spring constant in newtons per metre
\end{definition}
\begin{proof}
  This is the declared model definition of spring Constant In Newtons Per Meter.
\end{proof}
\begin{definition}[electromotive Force In Volts]
  \label{def:physics:phyx-mini-0979:phyxminiproblems-problemphyxmini0979-electromotiveforceinvolts}
  \lean{PhyXMiniProblems.ProblemPhyXMini0979.electromotiveForceInVolts}
  \uses{def:physics:phyx-mini-0979:phyxminiproblems-problemphyxmini0979-electromotiveforcemagnitude}
  Read an electromotive-force magnitude in volts
\end{definition}
\begin{proof}
  This is the declared model definition of electromotive Force In Volts.
\end{proof}
\begin{definition}[electric Current In Amperes]
  \label{def:physics:phyx-mini-0979:phyxminiproblems-problemphyxmini0979-electriccurrentinamperes}
  \lean{PhyXMiniProblems.ProblemPhyXMini0979.electricCurrentInAmperes}
  \uses{def:physics:phyx-mini-0979:phyxminiproblems-problemphyxmini0979-electriccurrentmagnitude}
  Read an electric-current magnitude in amperes
\end{definition}
\begin{proof}
  This is the declared model definition of electric Current In Amperes.
\end{proof}
\begin{definition}[force In Newtons]
  \label{def:physics:phyx-mini-0979:phyxminiproblems-problemphyxmini0979-forceinnewtons}
  \lean{PhyXMiniProblems.ProblemPhyXMini0979.forceInNewtons}
  \uses{def:physics:phyx-mini-0979:phyxminiproblems-problemphyxmini0979-forcemagnitude}
  Read a force magnitude in newtons
\end{definition}
\begin{proof}
  This is the declared model definition of force In Newtons.
\end{proof}
\begin{definition}[damping Coefficient In Kilograms Per Second]
  \label{def:physics:phyx-mini-0979:phyxminiproblems-problemphyxmini0979-dampingcoefficientinkilogramspersecond}
  \lean{PhyXMiniProblems.ProblemPhyXMini0979.dampingCoefficientInKilogramsPerSecond}
  \uses{def:physics:phyx-mini-0979:phyxminiproblems-problemphyxmini0979-dampingcoefficientmagnitude}
  Read a damping coefficient in kilograms per second
\end{definition}
\begin{proof}
  This is the declared model definition of damping Coefficient In Kilograms Per Second.
\end{proof}
\begin{definition}[Figure Object]
  \label{def:physics:phyx-mini-0979:phyxminiproblems-problemphyxmini0979-figureobject}
  \lean{PhyXMiniProblems.ProblemPhyXMini0979.FigureObject}
  Objects visibly present in image 979.png
\end{definition}
\begin{proof}
  This is the declared model definition of Figure Object.
\end{proof}
\begin{definition}[Figure Label]
  \label{def:physics:phyx-mini-0979:phyxminiproblems-problemphyxmini0979-figurelabel}
  \lean{PhyXMiniProblems.ProblemPhyXMini0979.FigureLabel}
  Literal symbolic labels visible in image 979.png
\end{definition}
\begin{proof}
  This is the declared model definition of Figure Label.
\end{proof}
\begin{definition}[Spatial Direction]
  \label{def:physics:phyx-mini-0979:phyxminiproblems-problemphyxmini0979-spatialdirection}
  \lean{PhyXMiniProblems.ProblemPhyXMini0979.SpatialDirection}
  Named directions needed for the top-view geometry
\end{definition}
\begin{proof}
  This is the declared model definition of Spatial Direction.
\end{proof}
\begin{definition}[Axis Orientation]
  \label{def:physics:phyx-mini-0979:phyxminiproblems-problemphyxmini0979-axisorientation}
  \lean{PhyXMiniProblems.ProblemPhyXMini0979.AxisOrientation}
  Orientation of an object or marked span in the figure
\end{definition}
\begin{proof}
  This is the declared model definition of Axis Orientation.
\end{proof}
\begin{definition}[Electromagnetic Spring Figure]
  \label{def:physics:phyx-mini-0979:phyxminiproblems-problemphyxmini0979-electromagneticspringfigure}
  \lean{PhyXMiniProblems.ProblemPhyXMini0979.ElectromagneticSpringFigure}
  \uses{def:physics:phyx-mini-0979:phyxminiproblems-problemphyxmini0979-figureobject, def:physics:phyx-mini-0979:phyxminiproblems-problemphyxmini0979-figurelabel, def:physics:phyx-mini-0979:phyxminiproblems-problemphyxmini0979-spatialdirection, def:physics:phyx-mini-0979:phyxminiproblems-problemphyxmini0979-axisorientation}
  Literal qualitative content of the primary raster. It records labels and geometry only and contains no amplitude readout
\end{definition}
\begin{proof}
  This is the declared model definition of Electromagnetic Spring Figure.
\end{proof}
\begin{definition}[Electromagnetic Spring Setup]
  \label{def:physics:phyx-mini-0979:phyxminiproblems-problemphyxmini0979-electromagneticspringsetup}
  \lean{PhyXMiniProblems.ProblemPhyXMini0979.ElectromagneticSpringSetup}
  \uses{def:physics:phyx-mini-0979:phyxminiproblems-problemphyxmini0979-massmagnitude, def:physics:phyx-mini-0979:phyxminiproblems-problemphyxmini0979-lengthmagnitude, def:physics:phyx-mini-0979:phyxminiproblems-problemphyxmini0979-timemagnitude, def:physics:phyx-mini-0979:phyxminiproblems-problemphyxmini0979-speedmagnitude, def:physics:phyx-mini-0979:phyxminiproblems-problemphyxmini0979-magneticfluxdensitymagnitude, def:physics:phyx-mini-0979:phyxminiproblems-problemphyxmini0979-resistancemagnitude, def:physics:phyx-mini-0979:phyxminiproblems-problemphyxmini0979-springconstantmagnitude, def:physics:phyx-mini-0979:phyxminiproblems-problemphyxmini0979-electromotiveforcemagnitude, def:physics:phyx-mini-0979:phyxminiproblems-problemphyxmini0979-electriccurrentmagnitude, def:physics:phyx-mini-0979:phyxminiproblems-problemphyxmini0979-forcemagnitude, def:physics:phyx-mini-0979:phyxminiproblems-problemphyxmini0979-dampingcoefficientmagnitude, def:physics:phyx-mini-0979:phyxminiproblems-problemphyxmini0979-spatialdirection, def:physics:phyx-mini-0979:phyxminiproblems-problemphyxmini0979-electromagneticspringfigure}
  Independent physical quantities and observables for the apparatus. The amplitude envelope is an observable indexed by clock time; it is not defined from the requested answer
\end{definition}
\begin{proof}
  This is the declared model definition of Electromagnetic Spring Setup.
\end{proof}
\begin{definition}[Matches Primary Electromagnetic Spring Figure]
  \label{def:physics:phyx-mini-0979:phyxminiproblems-problemphyxmini0979-matchesprimaryelectromagneticspringfigure}
  \lean{PhyXMiniProblems.ProblemPhyXMini0979.MatchesPrimaryElectromagneticSpringFigure}
  \uses{def:physics:phyx-mini-0979:phyxminiproblems-problemphyxmini0979-electromagneticspringsetup}
  Every qualitative label and relation directly read from image 979.png
\end{definition}
\begin{proof}
  This is the declared model definition of Matches Primary Electromagnetic Spring Figure.
\end{proof}
\begin{definition}[Matches Electromagnetic Spring Scenario]
  \label{def:physics:phyx-mini-0979:phyxminiproblems-problemphyxmini0979-matcheselectromagneticspringscenario}
  \lean{PhyXMiniProblems.ProblemPhyXMini0979.MatchesElectromagneticSpringScenario}
  \uses{def:physics:phyx-mini-0979:phyxminiproblems-problemphyxmini0979-electromagneticspringsetup}
  Qualitative prose assumptions about the idealized apparatus
\end{definition}
\begin{proof}
  This is the declared model definition of Matches Electromagnetic Spring Scenario.
\end{proof}
\begin{definition}[Electromagnetic Spring Problem Data]
  \label{def:physics:phyx-mini-0979:phyxminiproblems-problemphyxmini0979-electromagneticspringproblemdata}
  \lean{PhyXMiniProblems.ProblemPhyXMini0979.ElectromagneticSpringProblemData}
  \uses{def:physics:phyx-mini-0979:phyxminiproblems-problemphyxmini0979-massinkilograms, def:physics:phyx-mini-0979:phyxminiproblems-problemphyxmini0979-lengthinmeters, def:physics:phyx-mini-0979:phyxminiproblems-problemphyxmini0979-timeinseconds, def:physics:phyx-mini-0979:phyxminiproblems-problemphyxmini0979-speedinmeterspersecond, def:physics:phyx-mini-0979:phyxminiproblems-problemphyxmini0979-magneticfluxdensityinteslas, def:physics:phyx-mini-0979:phyxminiproblems-problemphyxmini0979-resistanceinohms, def:physics:phyx-mini-0979:phyxminiproblems-problemphyxmini0979-springconstantinnewtonspermeter, def:physics:phyx-mini-0979:phyxminiproblems-problemphyxmini0979-electromagneticspringsetup}
  The numerical values stated in the prose, expressed as exact coherent-SI readouts. 40.0 cm and 10.0 cm are converted to metres
\end{definition}
\begin{proof}
  This is the declared model definition of Electromagnetic Spring Problem Data.
\end{proof}
\begin{definition}[Electromagnetic Spring Laws]
  \label{def:physics:phyx-mini-0979:phyxminiproblems-problemphyxmini0979-electromagneticspringlaws}
  \lean{PhyXMiniProblems.ProblemPhyXMini0979.ElectromagneticSpringLaws}
  \uses{def:physics:phyx-mini-0979:phyxminiproblems-problemphyxmini0979-massinkilograms, def:physics:phyx-mini-0979:phyxminiproblems-problemphyxmini0979-lengthinmeters, def:physics:phyx-mini-0979:phyxminiproblems-problemphyxmini0979-timeinseconds, def:physics:phyx-mini-0979:phyxminiproblems-problemphyxmini0979-speedinmeterspersecond, def:physics:phyx-mini-0979:phyxminiproblems-problemphyxmini0979-magneticfluxdensityinteslas, def:physics:phyx-mini-0979:phyxminiproblems-problemphyxmini0979-resistanceinohms, def:physics:phyx-mini-0979:phyxminiproblems-problemphyxmini0979-springconstantinnewtonspermeter, def:physics:phyx-mini-0979:phyxminiproblems-problemphyxmini0979-electromotiveforceinvolts, def:physics:phyx-mini-0979:phyxminiproblems-problemphyxmini0979-electriccurrentinamperes, def:physics:phyx-mini-0979:phyxminiproblems-problemphyxmini0979-forceinnewtons, def:physics:phyx-mini-0979:phyxminiproblems-problemphyxmini0979-dampingcoefficientinkilogramspersecond, def:physics:phyx-mini-0979:phyxminiproblems-problemphyxmini0979-electromagneticspringsetup}
  Generic physical laws used to derive the damping and amplitude. None mentions the answer choice 5.13 cm or fixes the amplitude at the observation time
\end{definition}
\begin{proof}
  This is the declared model definition of Electromagnetic Spring Laws.
\end{proof}
\begin{theorem}[magnetic Damping Coefficient eq]
  \label{thm:physics:phyx-mini-0979:phyxminiproblems-problemphyxmini0979-magneticdampingcoefficient-eq}
  \lean{PhyXMiniProblems.ProblemPhyXMini0979.magneticDampingCoefficient_eq}
  \uses{def:physics:phyx-mini-0979:phyxminiproblems-problemphyxmini0979-lengthinmeters, def:physics:phyx-mini-0979:phyxminiproblems-problemphyxmini0979-magneticfluxdensityinteslas, def:physics:phyx-mini-0979:phyxminiproblems-problemphyxmini0979-resistanceinohms, def:physics:phyx-mini-0979:phyxminiproblems-problemphyxmini0979-dampingcoefficientinkilogramspersecond, def:physics:phyx-mini-0979:phyxminiproblems-problemphyxmini0979-electromagneticspringsetup, def:physics:phyx-mini-0979:phyxminiproblems-problemphyxmini0979-electromagneticspringproblemdata, def:physics:phyx-mini-0979:phyxminiproblems-problemphyxmini0979-electromagneticspringlaws}
  The induction laws make the magnetic drag coefficient B² W² / R
\end{theorem}
\begin{proof}
  The conclusion follows from the governing relations and figure readouts named in the statement of magnetic Damping Coefficient eq.
\end{proof}
\begin{theorem}[magnetic Damping Coefficient eq eight twenty fifths]
  \label{thm:physics:phyx-mini-0979:phyxminiproblems-problemphyxmini0979-magneticdampingcoefficient-eq-eight-twenty-fifths}
  \lean{PhyXMiniProblems.ProblemPhyXMini0979.magneticDampingCoefficient_eq_eight_twenty_fifths}
  \uses{def:physics:phyx-mini-0979:phyxminiproblems-problemphyxmini0979-dampingcoefficientinkilogramspersecond, def:physics:phyx-mini-0979:phyxminiproblems-problemphyxmini0979-electromagneticspringsetup, def:physics:phyx-mini-0979:phyxminiproblems-problemphyxmini0979-electromagneticspringproblemdata, def:physics:phyx-mini-0979:phyxminiproblems-problemphyxmini0979-electromagneticspringlaws}
  For the stated apparatus, the magnetic damping coefficient is 0.32 kg/s
\end{theorem}
\begin{proof}
  The conclusion follows from the governing relations and figure readouts named in the statement of magnetic Damping Coefficient eq eight twenty fifths.
\end{proof}
\begin{theorem}[oscillator is Underdamped]
  \label{thm:physics:phyx-mini-0979:phyxminiproblems-problemphyxmini0979-oscillator-isunderdamped}
  \lean{PhyXMiniProblems.ProblemPhyXMini0979.oscillator_isUnderdamped}
  \uses{def:physics:phyx-mini-0979:phyxminiproblems-problemphyxmini0979-electromagneticspringsetup, def:physics:phyx-mini-0979:phyxminiproblems-problemphyxmini0979-electromagneticspringproblemdata, def:physics:phyx-mini-0979:phyxminiproblems-problemphyxmini0979-electromagneticspringlaws}
  The calibrated oscillator lies in the oscillatory, underdamped regime
\end{theorem}
\begin{proof}
  The conclusion follows from the governing relations and figure readouts named in the statement of oscillator is Underdamped.
\end{proof}
\begin{theorem}[decay Rate eq two fifteenths]
  \label{thm:physics:phyx-mini-0979:phyxminiproblems-problemphyxmini0979-decayrate-eq-two-fifteenths}
  \lean{PhyXMiniProblems.ProblemPhyXMini0979.decayRate_eq_two_fifteenths}
  \uses{def:physics:phyx-mini-0979:phyxminiproblems-problemphyxmini0979-electromagneticspringsetup, def:physics:phyx-mini-0979:phyxminiproblems-problemphyxmini0979-electromagneticspringproblemdata, def:physics:phyx-mini-0979:phyxminiproblems-problemphyxmini0979-electromagneticspringlaws}
  The amplitude decay rate γ/(2M) is 2/15 s⁻¹
\end{theorem}
\begin{proof}
  The conclusion follows from the governing relations and figure readouts named in the statement of decay Rate eq two fifteenths.
\end{proof}
\begin{theorem}[amplitude at observation time eq]
  \label{thm:physics:phyx-mini-0979:phyxminiproblems-problemphyxmini0979-amplitude-at-observation-time-eq}
  \lean{PhyXMiniProblems.ProblemPhyXMini0979.amplitude_at_observation_time_eq}
  \uses{def:physics:phyx-mini-0979:phyxminiproblems-problemphyxmini0979-lengthinmeters, def:physics:phyx-mini-0979:phyxminiproblems-problemphyxmini0979-electromagneticspringsetup, def:physics:phyx-mini-0979:phyxminiproblems-problemphyxmini0979-electromagneticspringproblemdata, def:physics:phyx-mini-0979:phyxminiproblems-problemphyxmini0979-electromagneticspringlaws}
  At 5 s, the model gives the exact amplitude 0.1 exp (-2/3) metres
\end{theorem}
\begin{proof}
  The conclusion follows from the governing relations and figure readouts named in the statement of amplitude at observation time eq.
\end{proof}
\begin{definition}[Rounds To Hundredth Centimeter]
  \label{def:physics:phyx-mini-0979:phyxminiproblems-problemphyxmini0979-roundstohundredthcentimeter}
  \lean{PhyXMiniProblems.ProblemPhyXMini0979.RoundsToHundredthCentimeter}
  \uses{def:physics:phyx-mini-0979:phyxminiproblems-problemphyxmini0979-lengthmagnitude, def:physics:phyx-mini-0979:phyxminiproblems-problemphyxmini0979-lengthinmeters}
  RoundsToHundredthCentimeter length n means that the length, expressed in centimetres, lies within half of 0.01 cm of n / 100 centimetres
\end{definition}
\begin{proof}
  This is the declared model definition of Rounds To Hundredth Centimeter.
\end{proof}
% --- Archon declaration topology end ---
% --- Archon physics formalization source end ---
